\documentclass[serif]{beamer}
\input{sns_beamer_preamble.tex}

\begin{document}
\begin{frame}{평가원이 던지는 힌트}
    따끈따끈한 9월 모의평가 두 문제 보겠습니다.
    \begin{hint}[2027학년도 9월 모의평가 공통 10번]
        상수 $a$ $(a>1)$에 대하여 직선 $y=7$이 두 곡선 $y=16a^x$, $y=\dfrac{1}{4}a^x$와 만나는 점을 각각 $\mathrm A$, $\mathrm B$라 하자.
        $\overline{\mathrm{AB}}=4$일 때, $a$의 값은? [4점]
    \end{hint}
    \begin{hint}[2027학년도 9월 모의평가 미적분 30번]
        함수 $f(x)=2^{2x+\frac{1}{2}\cos\pi x}$의 역함수 $g(x)$는 양의 실수 전체의 집합에서 연속인 도함수를 갖는다.
        \[
            \frac{1}{4}\int_{f(1)}^{f(3)}g(x)\,dx+\int_{f(0)}^{f(2)}\frac{2^{4g(x)}g^\prime(x)}{x}\,dx=\frac{q}{p}\sqrt{2}
        \]
        일 때, $p+q$의 값을 구하시오. (단, $p$와 $q$는 서로소인 자연수이다.) [4점]
    \end{hint}
\end{frame}

\begin{frame}
    먼저 공통 10번을 보겠습니다. 두 곡선의 식을 살짝 바꾸면
    \[
        y=16a^x=a^{x+\log_a16}=a^{x+4\log_a2},\quad y=\frac{1}{4}a^x=a^{x+\log_a\frac{1}{4}}=a^{x-2\log_a2}
    \]
    이므로 평행이동을 생각하면 자연스럽게 $\overline{\mathrm{AB}}=6\log_a2$임을 알 수 있습니다. 따라서
    \[
        6\log_a2=4~\longrightarrow~\log_a2=\frac{2}{3}~\longrightarrow~\log_2a=\frac{3}{2}~\longrightarrow~a=2^{3/2}=2\sqrt{2}
    \]
    가 됩니다. 이 문항에서 지수함수의 $x$축 방향으로의 평행이동과 상수배가 동일함을 명심하고 미적분 30번을 보겠습니다. (이 문항에서 숫자 $7$은 장식이라는 사실도!)
\end{frame}

\begin{frame}{}
    미적분 30번입니다. 두 번째 항에서 치환적분 냄새를 맡았습니다. $g(x)=t$라 하면 $g^\prime(x)\,dx=dt$이므로
    \[
        [\text{2}]=\int_{f(0)}^{f(2)}\frac{2^{4g(x)}g^\prime(x)}{x}\,dx=\int_{0}^{2}\frac{2^{4t}}{f(t)}\,dt
    \]
    입니다. 먼저 $f(t)$를 대입하면
    \[
        [\text{2}]=\int_{0}^{2}\frac{2^{4t}}{f(t)}\,dt=\int_{0}^{2}2^{4t}\cdot 2^{-2t-\frac{1}{2}\cos\pi t}\,dt=\int_{0}^{2}2^{2t-\frac{1}{2}\cos\pi t}\,dt
    \]
    인데, 아마 여기서 막힐 겁니다. 그러니 다른 항도 살펴봅시다. 기출문제를 열심히 봤다면 \\ 다음 식을 떠올릴 수 있습니다.
    \[
        [\text{1}]=\frac{1}{4}\int_{f(1)}^{f(3)}g(x)\,dx=\frac{3f(3)-f(1)}{4}-\frac{1}{4}\int_{1}^{3}f(x)\,dx
    \]
    감을 잡으셨나요?
\end{frame}

\begin{frame}
    이제 머릿속에서
    \begin{enumerate}[-]
        \item $2t-\frac{1}{2}\cos\pi t$와 $f(x)=2x+\frac{1}{2}\cos\pi x$의 위화감
        \item 적분 구간 $[0,2]$와 $[1,3]$의 길이가 같다는 사실
        \item 삼각함수의 성질을 이용한 정적분 식 조작 [기출문제의 경험]
        \item 지수함수의 $x$축 방향 평행이동 [평가원의 힌트: 공통 10번]
    \end{enumerate}
    이 모든 것을 합쳐서 뭔가 적분 구간을 평행이동하고 싶다는 욕구가 들어야 합니다. \\
    여기서 마지막 항목은 평행이동으로 $\frac{1}{4}$을 얻는다는 확신을 얻기 위해 꼭 필요합니다.
    \begin{align*}
        [\text{2}]&=\int_{0}^{2}2^{2t-\frac{1}{2}\cos\pi t}\,dt\\[3pt]
        &=\int_{1}^{3}2^{2(t-1)-\frac{1}{2}\cos\pi(t-1)}\,dt\\[3pt]
        &=\int_{1}^{3}2^{2t-2+\frac{1}{2}\cos\pi t}\,dt\\[3pt]
        &=\frac{1}{4}\int_{1}^{3}f(t)\,dt
    \end{align*}
    임을 알 수 있고, 따라서
    \[
        (\text{준 식})=[\text{1}]+[\text{2}]=\frac{3f(3)-f(1)}{4}=\frac{47}{2}\sqrt{2}
    \]
    이므로 $p+q=47+2=49$입니다.
\end{frame}

\begin{frame}
    솔직히 저도 다음 기출문제가 없었다면 많이 헤맸을 겁니다.
    \begin{hint}[2025년 7월 학력평가 미적분 30번]
        함수 $\displaystyle f(x)=\int_0^xe^{\cos\pi t}\,dt$의 역함수를 $g(x)$라 할 때, 실수 전체의 집합에서 도함수가\\
        연속인 함수 $h(x)$가 모든 실수 $x$에 대하여
        \[
            h(g(x)+2)=2x^3+6f(1)x^2+1
        \]
        을 만족시킨다. $\displaystyle\int_3^7\frac{h^\prime(x)}{f(x)}\,dx=k\times\{f(1)\}^2$일 때, 실수 $k$의 값을 구하시오. [4점]
    \end{hint}
    기출문제 많이 푸세요.
\end{frame}
\end{document}